\documentclass[10pt,a4paper]{article}
\usepackage[a4paper,margin=0.7in]{geometry}
\usepackage{listings}
\usepackage{xcolor}
\usepackage{appendix}

% --- Verilog listing style ----------------------------------------------------
\definecolor{kw}{rgb}{0.10,0.20,0.65}
\definecolor{cm}{rgb}{0.30,0.55,0.30}
\definecolor{st}{rgb}{0.65,0.10,0.10}
\definecolor{bg}{rgb}{0.97,0.97,0.97}

\lstdefinelanguage{Verilog}{
  morekeywords={module,endmodule,input,output,inout,wire,reg,parameter,
    localparam,assign,always,posedge,negedge,begin,end,if,else,case,endcase,
    default,for,generate,endgenerate,genvar,signed,unsigned,initial,integer,
    function,endfunction,task,endtask,return,or,and,not,xor},
  sensitive=true,
  morecomment=[l]{//},
  morecomment=[s]{/*}{*/},
  morestring=[b]"
}

\lstset{
  language=Verilog,
  basicstyle=\ttfamily\scriptsize,
  keywordstyle=\color{kw}\bfseries,
  commentstyle=\color{cm}\itshape,
  stringstyle=\color{st},
  backgroundcolor=\color{bg},
  breaklines=true,
  showstringspaces=false,
  numbers=left,
  numberstyle=\tiny\color{gray},
  numbersep=4pt,
  xleftmargin=14pt,
  framexleftmargin=12pt,
  frame=single,
  rulecolor=\color{lightgray},
  tabsize=2,
  aboveskip=4pt,
  belowskip=4pt,
  captionpos=b
}

\setlength{\parindent}{0pt}
\setlength{\parskip}{2pt}

\begin{document}

\appendix
\section{Code Snippets of Key RTL Modules}
\label{appendix:rtl}

This appendix collects representative code excerpts from the Verilog
implementation of the GNN accelerator. Snippets focus on the
computational kernels (parallel MAC neuron, ReLU/saturation, layer
normalization, pipelined edge network), the top-level phase FSM, and
the small primitive blocks (counter, multiplier).

% =============================================================================
\subsection{Parallel MAC Neuron (\texttt{neuron.v})}
4-lane parallel multiply--accumulate; processes 4 weight$\times$data
pairs per cycle and accumulates with a 2-level adder tree. Cuts counter
cycles by~$4\times$ vs the scalar version (e.g. 192$\to$48 for L1 edge).

\begin{lstlisting}
module neuron #(parameter NEURON_WIDTH=8, DATA_BITS=16, WEIGHT_BITS=32,
                BIAS_BITS=16, WeightFile="...", BiasFile="...")(
    input clk, rstn, activation_function,
    input  signed [NEURON_WIDTH*DATA_BITS-1:0] data_in_flat,
    input  [31:0] counter,
    output signed [DATA_BITS+WEIGHT_BITS+8-1:0] data_out);

  localparam ACCUM_BITS = DATA_BITS+WEIGHT_BITS+8, N=4;
  reg signed [WEIGHT_BITS-1:0] weights [0:NEURON_WIDTH-1];
  reg signed [BIAS_BITS-1:0]   bias_mem [0:0];
  initial begin
    $readmemb(WeightFile, weights);
    $readmemb(BiasFile,   bias_mem);
  end

  // 4 parallel multiply lanes
  wire signed [DATA_BITS-1:0]   x [0:N-1];
  wire signed [WEIGHT_BITS-1:0] w [0:N-1];
  reg  signed [DATA_BITS+WEIGHT_BITS-1:0] prod [0:N-1];
  genvar k;
  generate for (k=0; k<N; k=k+1) begin : lane
    wire [31:0] idx = (counter>=1) ? (counter-1)*N+k : 0;
    assign w[k] = (counter>=1 && idx<NEURON_WIDTH) ? weights[idx] : 0;
    assign x[k] = (counter>=1 && idx<NEURON_WIDTH) ?
                  data_in_flat[idx*DATA_BITS +: DATA_BITS] : 0;
    always @(posedge clk or negedge rstn)
      if (!rstn) prod[k] <= 0; else prod[k] <= w[k]*x[k];
  end endgenerate

  // 2-level adder tree: (p0+p1) + (p2+p3)
  wire signed [DATA_BITS+WEIGHT_BITS:0] s01, s23;
  assign s01 = {{1{prod[0][DATA_BITS+WEIGHT_BITS-1]}},prod[0]}
             + {{1{prod[1][DATA_BITS+WEIGHT_BITS-1]}},prod[1]};
  assign s23 = {{1{prod[2][DATA_BITS+WEIGHT_BITS-1]}},prod[2]}
             + {{1{prod[3][DATA_BITS+WEIGHT_BITS-1]}},prod[3]};
  wire signed [DATA_BITS+WEIGHT_BITS+1:0] partial = s01 + s23;

  // Accumulator: starts at counter=2 (1cy mult latency)
  reg signed [ACCUM_BITS-1:0] accumulator, data_out_r;
  always @(posedge clk or negedge rstn) begin
    if (!rstn)            begin accumulator<=0; data_out_r<=0; end
    else if (counter==1)  begin accumulator<=0; data_out_r<=0; end
    else if (counter==2)  begin accumulator<=accumulator+partial;
                                data_out_r<=accumulator+partial; end
    else if (counter>=3)  begin accumulator<=accumulator+partial;
                                data_out_r<=accumulator; end
  end
  assign data_out = data_out_r;
endmodule
\end{lstlisting}

% =============================================================================
\subsection{ReLU + Saturation (\texttt{ReLu.v})}
Adds bias on the final counter tick, optionally applies ReLU, and
saturates the wide accumulator down to \texttt{DATA\_BITS} (Q4.4).

\begin{lstlisting}
module ReLu #(parameter DATA_BITS=8, WEIGHT_BITS=8, COUNTER_END=6,
              BIAS_BITS=8, ACCUM_EXTRA=8, WEIGHT_FRAC_BITS=6)(
    input clk, rstn, activation_function, input [31:0] counter,
    input signed [DATA_BITS+WEIGHT_BITS+ACCUM_EXTRA-1:0] mult_sum_in,
    input signed [BIAS_BITS-1:0] b,
    output reg signed [DATA_BITS-1:0] neuron_out);
  localparam signed [DATA_BITS-1:0] OUT_MAX = {1'b0,{(DATA_BITS-1){1'b1}}};
  localparam signed [DATA_BITS-1:0] OUT_MIN = {1'b1,{(DATA_BITS-1){1'b0}}};
  localparam signed [DATA_BITS+WEIGHT_BITS+ACCUM_EXTRA-1:0]
    MAX_VAL = {{(WEIGHT_BITS+ACCUM_EXTRA){1'b0}}, OUT_MAX},
    MIN_VAL = {{(WEIGHT_BITS+ACCUM_EXTRA){1'b1}}, OUT_MIN};
  reg signed [DATA_BITS+WEIGHT_BITS+ACCUM_EXTRA-1:0] biased_sum;
  reg signed [DATA_BITS-1:0] neuron_out_next;
  always @(posedge clk or negedge rstn) begin
    if (!rstn) neuron_out <= 0;
    else if (counter == COUNTER_END) begin
      biased_sum = mult_sum_in +
        {{(DATA_BITS+WEIGHT_BITS+ACCUM_EXTRA-BIAS_BITS){b[BIAS_BITS-1]}}, b};
      if (activation_function) begin           // ReLU + saturation
        if      (biased_sum < 0)        neuron_out_next = 0;
        else if (biased_sum > MAX_VAL)  neuron_out_next = OUT_MAX;
        else neuron_out_next =
          biased_sum[DATA_BITS+WEIGHT_BITS-WEIGHT_FRAC_BITS-1 -:DATA_BITS];
      end else begin                           // Linear + saturation
        if      (biased_sum > MAX_VAL) neuron_out_next = OUT_MAX;
        else if (biased_sum < MIN_VAL) neuron_out_next = OUT_MIN;
        else neuron_out_next =
          biased_sum[DATA_BITS+WEIGHT_BITS-WEIGHT_FRAC_BITS-1 -:DATA_BITS];
      end
      neuron_out <= neuron_out_next;
    end
  end
endmodule
\end{lstlisting}

% =============================================================================
\subsection{Top-Level Phase FSM (\texttt{topModule.v})}
Coordinates the pipeline phases. Edge and node encoders run in parallel
(shared no resources); message passing follows, then the edge decoder.

\begin{lstlisting}
localparam PHASE_IDLE=0, PHASE_EDGE_ENCODE=1, PHASE_EDGE_ENCODE_WAIT=2,
           PHASE_NODE_ENCODE=3, PHASE_NODE_ENCODE_WAIT=4,
           PHASE_MESSAGE_PASSING=5, PHASE_MESSAGE_PASSING_WAIT=6,
           PHASE_EDGE_DECODE=7, PHASE_EDGE_DECODE_WAIT=8, PHASE_DONE=9;

always @(posedge clk or negedge rstn) begin
  if (!rstn) phase_state <= PHASE_IDLE;
  else begin
    if (edge_done) edge_done_captured <= 1'b1;   // sticky latches
    if (node_done) node_done_captured <= 1'b1;
    case (phase_state)
      PHASE_IDLE: begin wait_counter<=10; phase_state<=PHASE_EDGE_ENCODE; end
      PHASE_EDGE_ENCODE:
        if (wait_counter>0) wait_counter<=wait_counter-1;
        else begin edge_start<=1; node_start<=1;        // start in parallel
                   phase_state<=PHASE_EDGE_ENCODE_WAIT; end
      PHASE_EDGE_ENCODE_WAIT:                           // wait BOTH done
        if ((edge_done|edge_done_captured) &&
            (node_done|node_done_captured)) begin
          edge_done_captured<=0; node_done_captured<=0;
          phase_state<=PHASE_MESSAGE_PASSING;
        end
      PHASE_MESSAGE_PASSING:
        if (wait_counter>0) wait_counter<=wait_counter-1;
        else begin mp_start<=1; phase_state<=PHASE_MESSAGE_PASSING_WAIT; end
      PHASE_MESSAGE_PASSING_WAIT:
        if (mp_done) phase_state<=PHASE_EDGE_DECODE;
      PHASE_EDGE_DECODE:
        if (wait_counter>0) wait_counter<=wait_counter-1;
        else begin dec_start<=1; phase_state<=PHASE_EDGE_DECODE_WAIT; end
      PHASE_EDGE_DECODE_WAIT: if (dec_done) phase_state<=PHASE_DONE;
      PHASE_DONE: phase_state<=PHASE_DONE;
    endcase
  end
end
assign processing_done = (phase_state == PHASE_DONE);
\end{lstlisting}

% =============================================================================
\subsection{Pipelined Edge Network (\texttt{Edge\_Network.v})}
3$\times$ L1 round-robin dispatch + L2/L3 pipeline via the
\texttt{CS\_L2\_GAP} state. Throughput becomes one layer time
(\texttt{max(L2,L3)}) instead of two.

\begin{lstlisting}
// 3x L1 instances dispatched round-robin from load FIFO
wire do_dispatch_a = !fifo_empty && (dispatch_rr==0) && !l1a_busy;
wire do_dispatch_b = !fifo_empty && (dispatch_rr==1) && !l1b_busy;
wire do_dispatch_c = !fifo_empty && (dispatch_rr==2) && !l1c_busy;

// L1->L2 inter-stage FIFO; L1 outputs collected into s_data[]
wire do_collect_a = Layer1a_valid && l1a_busy && !s_full;
wire do_collect_b = Layer1b_valid && l1b_busy && !s_full && !do_collect_a;
wire do_collect_c = Layer1c_valid && l1c_busy && !s_full
                                  && !do_collect_a && !do_collect_b;

// L2/L3 pipeline FSM
localparam CS_IDLE=0, CS_L2W=1, CS_L2_GAP=2, CS_L3W=3;
// CS_L2W   : L2 runs on edge k
// CS_L2_GAP: layer2_held=0 (counter resets), layer3_held=1 starts L3 on k;
//            if FIFO has edge k+1, load it -> L2 restarts next cycle
// CS_L3W   : L3 runs on edge k, L2 may run concurrently on edge k+1

// L2 ping-pong output buffer prevents L3 mid-run corruption
reg [DATA_BITS*EDGE_FEATURES-1:0] pp_buf [0:1];
reg pp_wr;
reg [DATA_BITS*EDGE_FEATURES-1:0] Layer2_out_r;  // stable L3 input
\end{lstlisting}

% =============================================================================
\subsection{Layer Normalization Top-Level FSM (\texttt{layer\_norm.v})}
Sequences the four sub-modules implementing
$y_i=((x_i-\mu)/\sqrt{\sigma^2+\epsilon})\gamma_i+\beta_i$ then optional
ReLU. Sub-module latencies: mean 6, var 7, inv\_sqrt 4, normalize 3,
relu 1.

\begin{lstlisting}
localparam [3:0] S_IDLE=0, S_CAPTURE=1, S_WAIT_MEAN=2, S_START_VAR=3,
  S_WAIT_VAR=4, S_START_SQRT=5, S_WAIT_SQRT=6, S_START_NORM=7,
  S_WAIT_NORM=8, S_START_RELU=9, S_WAIT_RELU=10, S_OUTPUT=11;
always @(*) begin
  next_state = state;
  case (state)
    S_IDLE       : if (valid_in) next_state = S_CAPTURE;
    S_CAPTURE    : next_state = S_WAIT_MEAN;
    S_WAIT_MEAN  : if (mean_valid)     next_state = S_START_VAR;
    S_START_VAR  : next_state = S_WAIT_VAR;
    S_WAIT_VAR   : if (var_valid)      next_state = S_START_SQRT;
    S_START_SQRT : next_state = S_WAIT_SQRT;
    S_WAIT_SQRT  : if (sqrt_valid)     next_state = S_START_NORM;
    S_START_NORM : next_state = S_WAIT_NORM;
    S_WAIT_NORM  : if (norm_valid)     next_state = S_START_RELU;
    S_START_RELU : next_state = S_WAIT_RELU;
    S_WAIT_RELU  : if (relu_valid)     next_state = S_OUTPUT;
    S_OUTPUT     : next_state = S_IDLE;
  endcase
end
\end{lstlisting}

% =============================================================================
\subsection{Counter \& Multiplier Primitives}
\begin{lstlisting}
module counter #(parameter END_COUNTER = 32'd0)(
    input clk, rstn, output reg [31:0] counter_out,
    output reg counter_donestatus);
  always @(posedge clk) begin
    if (!rstn) begin counter_out<=0; counter_donestatus<=0; end
    else if (counter_out >= END_COUNTER+1) begin
      counter_out <= END_COUNTER+1; counter_donestatus <= 1'b1;
    end else begin
      counter_out <= counter_out + 1'b1; counter_donestatus <= 1'b0;
    end
  end
endmodule

module multiplier #(parameter DATA_BITS=8, WEIGHT_BITS=8)(
    input clk, rstn,
    input  signed [WEIGHT_BITS-1:0] w,
    input  signed [DATA_BITS-1:0]   x,
    output reg signed [DATA_BITS+WEIGHT_BITS-1:0] mult_result);
  always @(posedge clk or negedge rstn)
    if (!rstn) mult_result <= 0;
    else       mult_result <= w * x;
endmodule
\end{lstlisting}

\end{document}

\section{Reset Strategy}

All RTL modules use a synchronous active-low reset (\texttt{rstn}).
Storage and BRAM-backed buffers reset their address counters but
preserve memory contents (which are loaded via \texttt{\$readmem*}
at elaboration time). The burst sequencer is the only stage with an
asynchronous-style \texttt{rst} input, retained for ease of testbench
control.

\section{Decoder}

The edge decoder mirrors the encoder structure but with a single
output dimension. \texttt{Edge\_Output\_Transform.v} flattens the
per-edge feature vector to a scalar prediction stream, which is then
written back to the storage subsystem for host readout.

\section{Verification flow}

\begin{enumerate}
  \item \texttt{gen\_expected.py} produces golden vectors in
        \texttt{mem\_files/}.
  \item iverilog runs the top-level testbench, reading the same
        \texttt{mem\_files/} via the storage init paths.
  \item \texttt{test.py} compares the per-edge output stream to the
        golden vectors, tolerating $\pm 1$ LSB of fixed-point error.
\end{enumerate}

\section{Parameter Quantization}

All hidden activations are stored in Q1.7. Weights are Q2.6 with
saturating clamps at the MAC output. Layer-norm gamma / beta are
stored in Q1.7. The inverse-square-root LUT is Q1.11. The choice
was driven by the smallest representable variance after the
8-feature aggregation step in the first MP block.

\section{Status}

The pipeline currently supports graphs up to 16 nodes / 32 edges
with feature dim 32. The decoder produces a per-edge scalar
prediction. End-to-end latency is approximately 3.5k cycles per
graph at 200 MHz.
